%! Author = lili
%! Date = 5/21/26

% S. 17 Bis Skript S. 22

In the ``Sequence Analysis 1'' class, we have already discussed a variant of the Universal
Alignment Algorithm that can be used to find a best (highest score) alignment of a (short)
pattern x within a (long) text y: semi-global alignment.
Often, we are interested in all matches with at least a certain score.
A simple modification makes this possible: We just
need to visit all predecessors of the final cell vE and check whether the corresponding
solutions fulfill the given criterion.
(In the case of semi-global alignment, these cells
correspond to the last row of the alignment matrix.)

In this section, we consider the cost-based formulation and therefore alignment matrix $D$.
We assume that a sensible cost function cost: $\mathcal{A} \rightarrow \rz_0^+$ for alignment columns is given
where $cost((\binom{a}{a})) = 0$, $cost((\binom{a}{b})) > 0$ for $a \neq b$ and all indels have the same cost $\gamma > 0$. \\
Then we consider the following problem variant:

\defin{approximate-string-matching-problem}

In fact, it suffices to discover the ending positions j of the substrings y' because we can
always reconstruct the whole substring and its starting position by backtracing.

\deffin{k_approximate_matches}

\section{Sellers' Algorithm}\label{sec:sellers'-algorithm}
For a change in style, here we present an explicit formulation of an algorithm that solves
the problem not as a recurrence, but in pseudo-code notation.
Algorithm 1 is known as
Sellers’ algorithm (Sellers, 1980).
Of course, it corresponds to the Universal Alignment
Algorithm variant discussed in the “Sequence Analysis 1” class for semi-global alignment,
except that we do not look at the final vertex vE to find the best match, but look at all
the vertices in the last row to identify all matches with $D(m, j) \leq k for 0 \leq j \leq n$.

Looking closely at Algorithm~\ref{alg:sellers}, we see that each iteration $j = 1, \dots, n$ transforms the
previous column $C_{j-1} \coloneqq D(\cdot, j - 1)$ of the edit matrix into the current column $C_j \coloneqq
D(\cdot, j)$, based on character $y[j]$.
Starting with $C_0$, Sellers’ algorithm effectively consists of
repeatedly computing $C_j$ from $C_j_{-1}$ and the next text character $y[j]$, for each $j = 1, \dots, n$.

\begin{algorithm}[h]
    \begin{algorithmic}
        todo
    \end{algorithmic}\label{alg:sellers}
\end{algorithm}

Some remarks about Sellers’ algorithm:

\paragraph{Since we only report end positions and costs,} there is no need to store back pointers
and the entire alignment matrix D in Algorithm~\ref{alg:sellers}.
The current and the previous
column suffice.
This decreases the space requirement to O(m), where m is the (short)
pattern length.
We still need to store the text, which is O(n), however this need not
necessarily be in memory.

\paragraph{If we have a very good match} with cost $< k$ at position $j$, the neighboring positions will also match with cost $\leq k$.
To avoid this redundancy, it makes sense to
post-process the output and only look for runs of local minima.
Without formally defining it, if $k = 3$ and the respective costs at positions $j - 2, \dots, j + 4$ are
$(3, 2, 1, 1, 2, 2, 3)$, the original algorithm will report all of these positions.
In most
applications, however, we are only interested in the locally minimal value 1, which
occurs as a run of length 2.
So it would make sense to report the run interval and
the minimum cost as $([j, j + 1], 1)$.

\begin{notte}
    Note that at no time in the algorithm we need to know the exact value of $C_j (m)$ as
    long as we can be sure that it exceeds k and therefore will not be reported.
\end{notte}

\section{Ukkonen’s Cutoff Algorithm}\label{sec:ukkonens-cutoff-algorithm}
The last of the above three observations leads to a considerable improvement of Sellers’
algorithm: If we know that in column $j$ all values after a certain index $\lei_j$ exceed $k$, we do
not need to compute them (e.g.\ we could assume that they are all equal to $\infty$ or $k + 1$).
The following definition captures this formally.

\deffin{last essential index}

The last essential index $\lei_j$ of column $C_j$ is $\lei_j \coloneqq \lei (C_j ) = \max\{i | C_j (i) \leq k\}$.
From
the definition it is clear that we cannot miss any k-approximate matches if we do not
compute the cells below the last essential index.

To set this in operation, the last essential index $\lei_0$ of the 0-th column is easily found,
because the scores increase monotonically from 0 to $m \cdot \gamma$.
The other columns, however,
are not necessarily monotone!
Therefore, as we move column-by-column through the
matrix, we have to make sure that we can efficiently find the last essential index of the
next column, given the previous column.
Consider the effects of the $C_{j-1}$-entries below
the last essential index on the next column $C_j$.
Since $C_{j-1} (i) > k$ for all $i > \lei_j-1$ and
costs are non-negative, these cells provide candidate values $> k for C_j (i)$ for $i > \lei_{j-1}+ 1$
which therefore do not need to be considered during the minimization.
The only chance
that such $C_j (i)$ remain bounded by $k$ is vertically, i.e., if $C_j (i) = C_j (i - 1) + \gamma$.
As soon
as $C_j (i) > k$ for some $i > \lei_{j-1}+ 1$, we know that the last essential index $\lei_j$ of $C_j$ occurs
before that $i$.
Algorithm~\ref{alg:ukkonen} shows the details.

\begin{algorithm}[h]
    \begin{algorithmic}
        todo
    \end{algorithmic}\label{alg:ukkonen}
\end{algorithm}

\begin{bsp}
    TODO S 19
\end{bsp}

\section{Search Schemes and Bidirectional Indices}\label{sec:search-schemes-and-bidirectional-indices}